Esta sección presenta la evaluación experimental del mecanismo de almacenamiento comprimido
integrado en el simulador. La comparación se realizó entre tres estrategias de representación
del estado cuántico: almacenamiento no comprimido, compresión sin pérdida mediante Blosc y
compresión con pérdida controlada mediante ZFP. El análisis se organizó sobre dos familias de
circuitos con perfiles de acceso contrastantes al vector de estado: búsqueda de Grover y
transformada cuántica de Fourier (QFT).

La campaña experimental cubrió 22, 23 y 24 qubits. En Grover se evaluaron instancias con
objetivo único y multiobjetivo; en QFT se estudiaron una superposición periódica y una
superposición de alta entropía con fases aleatorias. Esta combinación permite contrastar
casos con alta regularidad estructural frente a entradas donde la compresión dispone de poca
redundancia explotable.

Las tablas de resultados reportan memoria residente máxima, ahorro de memoria respecto a la
referencia no comprimida, tiempo total de ejecución y tiempo relativo expresado en múltiplos
de la referencia. En Blosc, al tratarse de una estrategia sin pérdida, el error máximo
absoluto frente a la referencia es cero. Este resultado se obtuvo comparando amplitud por
amplitud el vector completo exportado por la ejecución densa y el vector completo exportado
por la ejecución con Blosc sobre la misma carga de trabajo. En ZFP, la columna de error
reporta el máximo error absoluto medido mediante la misma comparación completa frente a la
referencia no comprimida. En la comparación complementaria con la
variante optimizada de Grover, en cambio, la lectura se concentra en memoria y tiempo, ya que
el propósito de esa subsección es aislar el efecto de comprimir un esquema cuya ventaja
principal proviene de la reducción analítica del estado.

En el caso de Grover se añade una comparación complementaria con la variante de
parametrización reducida descrita en la Sección~\ref{subsec:grover_reducido_impl}. Su
propósito es delimitar cuánto puede reducirse el costo computacional cuando la simulación
explota directamente la simetría global del algoritmo, más allá del efecto atribuible al
esquema de almacenamiento.
